% Category A; sources are intentionally not shown on slides.
\categorycover{HBF System Characterization\\and KV-Centric Serving}{When does HBF improve LLM serving, and when does it become a bottleneck?}{Characterizing HBF; HBF Sucks?; Hot--Cold Tiering}
% Source: https://arxiv.org/abs/2609.39131v2
\begin{frame}[t]{Characterizing High Bandwidth Flash for LLM Serving}
\vspace{.02cm}
\begin{minipage}[t][3.25cm][t]{\linewidth}
\begin{summarysection}{Motivation}
\item \key{Capacity pressure:} Long contexts and agentic calls compete for HBM space for weights, active KV and reusable prefixes.
\item \key{Retention matters:} More HBF capacity helps only when placement and scheduling account for read costs, slow writes and finite endurance.
\item \key{Outcome:} Trace-driven simulation improves completion time and estimated write lifetime. Modeled energy savings exclude static power.
\end{summarysection}
\end{minipage}
\vspace{.18cm}
\begin{summarysection}{Method}
\item \key{Organize:} Model an HBM--HBF--host hierarchy and compare weight placement and KV residency under different memory organizations.
\item \key{Retain:} Track prefix reuse and cache movement. Buffer growing KV in HBM and preserve useful prefixes to avoid reloads or recomputation.
\item \key{Schedule:} Prioritize cache hits and limit admission with buffered cache-aware scheduling, balancing batch growth against future KV reuse.
\end{summarysection}
\end{frame}
% Source: https://arxiv.org/abs/2608.11668v4
\begin{frame}[t]{HBF Sucks? A Full-Stack Characterization of\\High-Bandwidth Flash for KV-Centric LLM Serving}
\vspace{.02cm}
\begin{minipage}[t][3.25cm][t]{\linewidth}
\begin{summarysection}{Motivation}
\item \key{Drop-in risk:} Faster flash does not ensure faster serving when HBF replaces SSD storage but also consumes near-tier package resources.
\item \key{Traffic mismatch:} Hot KV remains in the near tier, leaving HBF a cold, write-heavy stream with limited useful reads to repay writes.
\item \key{Outcome:} The modeled Mooncake-style substitution can worsen latency and goodput. This is a conditional deployment result, not a rejection of HBF.
\end{summarysection}
\end{minipage}
\vspace{.18cm}
\begin{summarysection}{Method}
\item \key{Replay:} Extend TokenSim with GPU/near-tier/HBF costs and replay production KV-offload traces using capacity-matched SSD baselines.
\item \key{Diagnose:} Test three conditions: reads on the critical path, enough useful reads per write, and bandwidth that can be sustained.
\item \key{Constrain:} Combine thermal throttling and write-endurance models to distinguish peak device bandwidth from practical serving benefits.
\end{summarysection}
\end{frame}
% Source: https://arxiv.org/abs/2609.25782v1
\begin{frame}[t]{Hot--Cold Tiering of HBM and High Bandwidth Flash\\for Agentic LLM Serving}
\vspace{.02cm}
\begin{minipage}[t][3.25cm][t]{\linewidth}
\begin{summarysection}{Motivation}
\item \key{Idle state:} Agentic sessions pause for tools or users while retaining large KV histories that exceed HBM capacity.
\item \key{Two access patterns:} Active KV is read every decode step; paused KV is read on resume. Serving all KV from HBF wastes read energy.
\item \key{Outcome:} Trace-based evaluation expands retained-session capacity while keeping active decoding in HBM and reducing flash read power.
\end{summarysection}
\end{minipage}
\vspace{.18cm}
\begin{summarysection}{Method}
\item \key{Classify:} Use session activity to separate a small hot KV working set from a large cold pool of paused-session state.
\item \key{Place:} Keep active KV in HBM. Evict cold KV to HBF under capacity pressure; immutable blocks with valid HBF copies need no rewrite.
\item \key{Resume:} Prefetch paused KV back into HBM through buffered inter-tier transfers before execution resumes, overlapping movement with decode.
\end{summarysection}
\end{frame}
